\documentclass[../../main.tex]{subfiles}
\graphicspath{{\subfix{../../image/}}}

\begin{document}

\section{特征值与特征函数}\label{section:ch03-特征线}

我们在本节考虑边值问题
\begin{align}
\begin{cases}
Lw=\lambda w & \text{于}\ U,\\
w=0 & \text{于}\ \partial U,
\end{cases}
\end{align}
其中 $U$ 开、有界；回忆 $\lambda$ 是 $L$ 的特征值，只要式~\eqref{align:ch06-更高边界问题} 存在非平凡解 $w$。由\cref{section:ch06-弱解存在性} 发展的理论，我们回忆 $L$ 的特征值集合 $\Sigma$ 至多可数。

下面\cref{subsection:ch06-对称算子特征值} 中的定理是线性代数标准断言（实对称矩阵有实特征值及特征向量正交基）对椭圆 PDE 的类比。类似地，\cref{subsection:ch06-非对称算子特征值} 中的结果是 Perron--Frobenius 定理（正矩阵有实正特征值及相应正分量特征向量）的 PDE 版本（见 Gantmacher [Ga]）。

\subsection{对称椭圆算子的特征值}\label{subsection:ch06-对称算子特征值}

为简单起见，我们现考虑具有散度型的椭圆算子
\begin{align}
Lu=-\sum_{i,j=1}^{n}(a_{ij}u_{x_{i}})_{x_{j}},
\end{align}
其中 $a_{ij}\in C^{\infty}(\overline{U})$（$i,j=1,\cdots,n$）。我们假设通常的一致椭圆性条件成立，且如通常那样假设
\begin{align}
a_{ij}=a_{ji}\qquad (i,j=1,\cdots,n).
\end{align}
算子 $L$ 因而形式对称，特别地相伴双线性型 $B[,]$ 满足 $B[u,v]=B[v,u]$（$u,v\in H^{1}_{0}(U)$）。另设 $U$ 连通。

\begin{theorem}[对称椭圆算子的特征值]\label{theorem:ch06-对称算子特征值}
\begin{enumerate}[(i)]
\item $L$ 的每个特征值为实数。
\item 此外，若我们按（有限）重数重复每个特征值，则
\begin{align*}
\Sigma=\{\lambda_{k}\}_{k=1}^{\infty},
\end{align*}
其中
\begin{align*}
0<\lambda_{1}\leqslant\lambda_{2}\leqslant\lambda_{3}\leqslant\cdots
\end{align*}
且当 $k\to\infty$ 时 $\lambda_{k}\to\infty$。
\item 最后，存在 $L^{2}(U)$ 的一组正交基 $\{w_{k}\}_{k=1}^{\infty}$，其中 $w_{k}\in H^{1}_{0}(U)$ 为对应于 $\lambda_{k}$ 的特征函数：
\begin{align}
\begin{cases}
Lw_{k}=\lambda_{k}w_{k} & \text{于}\ U,\\
w_{k}=0 & \text{于}\ \partial U,
\end{cases}
\end{align}
$k=1,2,\cdots$。由于\cref{section:ch06-正则性} 的正则性理论，$w_{k}\in C^{\infty}(U)$；此外若 $\partial U$ 光滑，$w_{k}\in C^{\infty}(\overline{U})$，$k=1,2,\cdots$。
\end{enumerate}
\end{theorem}

\begin{proof}
1. 如\cref{section:ch06-弱解存在性} 中那样，
\begin{align*}
S\triangleq L^{-1}
\end{align*}
是把 $L^{2}(U)$ 映入自身的有界、线性、紧算子。

2. 我们进一步断言 $S$ 对称。为此，选 $f,g\in L^{2}(U)$。则 $Sf=u$ 意味着 $u\in H^{1}_{0}(U)$ 是
\begin{align*}
\begin{cases}
Lu=f & \text{于}\ U,\\
u=0 & \text{于}\ \partial U
\end{cases}
\end{align*}
的弱解；类似地 $Sg=v$ 意味着 $v\in H^{1}_{0}(U)$ 在弱意义下求解
\begin{align*}
\begin{cases}
Lv=g & \text{于}\ U,\\
v=0 & \text{于}\ \partial U.
\end{cases}
\end{align*}
因而
\begin{align*}
(Sf,g)=(u,g)=B[v,u],\qquad (f,Sg)=(f,v)=B[u,v].
\end{align*}
由于 $B[u,v]=B[v,u]$，我们看到对一切 $f,g\in L^{2}(U)$ 有 $(Sf,g)=(f,Sg)$。因此 $S$ 对称。

3. 另注意
\begin{align*}
(Sf,f)=(u,f)=B[u,u]\geqslant 0\qquad (f\in L^{2}(U)).
\end{align*}
因而\cref{section:appD-对称算子} 的紧对称算子理论蕴含：$S$ 的所有特征值为实、正，且有相应特征函数组成 $L^{2}(U)$ 的一组正交基。但另注意对 $\eta\neq 0$，$Sw=\eta w$ 当且仅当 $Lw=\lambda w$（$\lambda=\dfrac{1}{\eta}$）。定理随之成立。
\end{proof}

\textbf{Weyl 律。} 椭圆算子特征值分布的研究在数学物理中极为重要。关于特征值渐近分布的一个里程碑式断言由 H. Weyl 给出。对光滑、有界开集 $U\subset\mathbb{R}^{n}$ 中带零边界条件的Laplace算子的特征值 $\{\lambda_{k}\}_{k=1}^{\infty}$ 的特殊情形，Weyl 律断言
\begin{align}
\lim_{k\to\infty}\frac{\lambda_{k}}{k^{2/n}}=\frac{(2\pi)^{n}}{|U|\alpha(n)}.
\end{align}
这里 $|U|$ 表示 $U$ 的体积。

下面我们更仔细地考察 $L$ 的第一特征值。

\begin{definition}\label{definition:ch06-主特征值}
我们称 $\lambda_{1}>0$ 为 $L$ 的主特征值。
\end{definition}

\begin{theorem}[主特征值的变分原理]\label{theorem:ch06-主特征值变分}
\begin{enumerate}[(i)]
\item 我们有
\begin{align}
\lambda_{1}=\min\{B[u,u]\mid u\in H^{1}_{0}(U),\ \|u\|_{L^{2}}=1\}.
\end{align}
\item 此外，上述最小值由一函数 $w_{1}$ 达到，$w_{1}$ 在 $U$ 内为正，并求解
\begin{align*}
\begin{cases}
Lw_{1}=\lambda_{1}w_{1} & \text{于}\ U,\\
w_{1}=0 & \text{于}\ \partial U.
\end{cases}
\end{align*}
\item 最后，若 $u\in H^{1}_{0}(U)$ 是
\begin{align*}
\begin{cases}
Lu=\lambda_{1}u & \text{于}\ U,\\
u=0 & \text{于}\ \partial U
\end{cases}
\end{align*}
的任一弱解，则 $u$ 是 $w_{1}$ 的倍数。
\end{enumerate}
\end{theorem}

\begin{remark}\label{remark:ch06-主特征值}
\leavevmode
\begin{enumerate}[(i)]
\item 断言 (iii) 说主特征值 $\lambda_{1}$ 单重。特别地
\begin{align*}
0<\lambda_{1}<\lambda_{2}\leqslant\lambda_{3}\leqslant\cdots.
\end{align*}
\item 式~\eqref{align:ch06-Rayleigh} 是 Rayleigh 公式，等价于陈述
\begin{align*}
\lambda_{1}=\min_{\substack{u\in H^{1}_{0}(U)\\ u\not\equiv 0}}\frac{B[u,u]}{\|u\|_{L^{2}(U)}^{2}}.
\end{align*}
\end{enumerate}
\end{remark}

\begin{proof}
1. 由式~\eqref{align:ch06-更高边界问题}, 我们看到
\begin{align}
B[w_{k},w_{k}]=\lambda_{k}\|w_{k}\|_{L^{2}(U)}^{2}=\lambda_{k}
\end{align}
且
\begin{align}
B[w_{k},w_{l}]=\lambda_{k}(w_{k},w_{l})=0
\end{align}
对 $k,l=1,2,\cdots$、$k\neq l$。

2. 由于 $\{w_{k}\}_{k=1}^{\infty}$ 是 $L^{2}(U)$ 的正交基，若 $u\in H^{1}_{0}(U)$、$\|u\|_{L^{2}(U)}=1$，我们可写
\begin{align}
u=\sum_{k=1}^{\infty}d_{k}w_{k}
\end{align}
其中 $d_{k}=(u,w_{k})_{L^{2}(U)}$，级数在 $L^{2}(U)$ 中收敛（见\cref{section:appD-Hilbert空间}）。此外
\begin{align}
\sum_{k=1}^{\infty}d_{k}^{2}=\|u\|_{L^{2}(U)}^{2}=1.
\end{align}

3. 此外由前述两式，我们看到 $\left\{\dfrac{w_{k}}{\lambda_{k}^{1/2}}\right\}_{k=1}^{\infty}$ 是 $H^{1}_{0}(U)$ 赋予新内积 $B[,]$ 后的正交子集。我们进一步断言 $\left\{\dfrac{w_{k}}{\lambda_{k}^{1/2}}\right\}_{k=1}^{\infty}$ 事实上是 $H^{1}_{0}(U)$ 赋予此新内积后的一组正交基。为此，只需验证
\begin{align*}
B[w_{k},u]=0\qquad (k=1,2,\cdots)
\end{align*}
蕴含 $u\equiv 0$。但此断言显然成立，因为恒等式
\begin{align*}
B[w_{k},u]=\lambda_{k}(w_{k},u)=0\qquad (k=1,\cdots)
\end{align*}
迫使 $u\equiv 0$，因为 $\{w_{k}\}_{k=1}^{\infty}$ 是 $L^{2}(U)$ 的基。因而
\begin{align*}
u=\sum_{k=1}^{\infty}\mu_{k}\frac{w_{k}}{\lambda_{k}^{1/2}}
\end{align*}
其中 $\mu_{k}=B\left[u,\dfrac{w_{k}}{\lambda_{k}^{1/2}}\right]$，级数在 $H^{1}_{0}(U)$ 中收敛。但由式~\eqref{align:ch06-更高边界问题}, $\mu_{k}=d_{k}\lambda_{k}^{1/2}$；因而级数~\eqref{align:ch06-更高边界问题} 事实上也在 $H^{1}_{0}(U)$ 中收敛。

4. 于是前述各式蕴含
\begin{align*}
B[u,u]=\sum_{k=1}^{\infty}d_{k}^{2}\lambda_{k}\geqslant\lambda_{1}\quad\text{由式~\eqref{align:ch06-正交}}.
\end{align*}
由于对 $u=w_{1}$ 等号成立，我们得到式~\eqref{align:ch06-Rayleigh}。

5. 我们 next 断言：若 $u\in H^{1}_{0}(U)$、$\|u\|_{L^{2}(U)}=1$，则 $u$ 是
\begin{align}
\begin{cases}
Lu=\lambda_{1}u & \text{于}\ U,\\
u=0 & \text{于}\ \partial U
\end{cases}
\end{align}
的弱解当且仅当
\begin{align}
B[u,u]=\lambda_{1}.
\end{align}
显然式~\eqref{align:ch06-正交} 蕴含式~\eqref{align:ch06-正交}。反之，设式~\eqref{align:ch06-正交} 成立。则如上写 $d_{k}=(u,w_{k})$，有
\begin{align}
\sum_{k=1}^{\infty}d_{k}^{2}\lambda_{1}=\lambda_{1}=B[u,u]=\sum_{k=1}^{\infty}d_{k}^{2}\lambda_{k}.
\end{align}
因而
\begin{align}
\sum_{k=1}^{\infty}(\lambda_{k}-\lambda_{1})d_{k}^{2}=0.
\end{align}
于是若 $\lambda_{k}>\lambda_{1}$，则 $d_{k}=(u,w_{k})=0$。由于 $\lambda_{1}$ 重数有限，得
\begin{align}
u=\sum_{k=1}^{m}(u,w_{k})w_{k}
\end{align}
某个 $m$，其中 $Lw_{k}=\lambda_{1}w_{k}$。因此
\begin{align}
Lu=\sum_{k=1}^{m}(u,w_{k})Lw_{k}=\lambda_{1}u.
\end{align}
这就证明了式~\eqref{align:ch06-正交}。

6. 下面我们将证明：若 $u\in H^{1}_{0}(U)$ 是式~\eqref{align:ch06-正交} 的弱解、$u\not\equiv 0$，则或者
\begin{align}
u>0\quad\text{于}\ U
\end{align}
或者
\begin{align}
u<0\quad\text{于}\ U.
\end{align}
为此，我们不失一般性设 $\|u\|_{L^{2}}=1$，注意
\begin{align}
\alpha+\beta=1
\end{align}
其中
\begin{align*}
\alpha\triangleq\int_{U}(u^{+})^{2}\,\mathrm{d}x,\qquad
\beta\triangleq\int_{U}(u^{-})^{2}\,\mathrm{d}x.
\end{align*}
此外由于 $u^{\pm}\in H^{1}_{0}(U)$，且
\begin{align*}
Du^{+}=
\begin{cases}
Du, & \text{a.e. 于}\ \{u\geqslant 0\},\\
0, & \text{a.e. 于}\ \{u\leqslant 0\},
\end{cases}
\qquad
Du^{-}=
\begin{cases}
0, & \text{a.e. 于}\ \{u\geqslant 0\},\\
-Du, & \text{a.e. 于}\ \{u\leqslant 0\}
\end{cases}
\end{align*}
（见第 5 章习题 18），我们有 $B[u^{+},u^{-}]=0$。因而
\begin{align*}
\lambda_{1}=B[u,u]=B[u^{+},u^{+}]+B[u^{-},u^{-}]
\geqslant\lambda_{1}\|u^{+}\|_{L^{2}(U)}^{2}+\lambda_{1}\|u^{-}\|_{L^{2}(U)}^{2}\quad\text{由式~\eqref{align:ch06-Rayleigh}}
=(\alpha+\beta)\lambda_{1}=\lambda_{1}.
\end{align*}
但于是我们看到上述不等式事实上必为等式，故
\begin{align*}
B[u^{+},u^{+}]=\lambda_{1}\|u^{+}\|_{L^{2}(U)}^{2},\qquad
B[u^{-},u^{-}]=\lambda_{1}\|u^{-}\|_{L^{2}(U)}^{2}.
\end{align*}
因此第 5 步中证明的断言断言
\begin{align}
\begin{cases}
Lu^{+}=\lambda_{1}u^{+} & \text{于}\ U,\\
u^{+}=0 & \text{于}\ \partial U
\end{cases}
\end{align}
且
\begin{align}
\begin{cases}
Lu^{-}=\lambda_{1}u^{-} & \text{于}\ U,\\
u^{-}=0 & \text{于}\ \partial U
\end{cases}
\end{align}
在弱意义下。

7. 其次，由于系数 $a_{ij}$ 光滑，我们由式~\eqref{align:ch06-u+方程} 推得 $u^{+}\in C^{\infty}(U)$ 且
\begin{align*}
Lu^{+}=\lambda_{1}u^{+}\geqslant 0\quad\text{于}\ U.
\end{align*}
函数 $u^{+}$ 因而为上解。于是强极值原理蕴含或者 $u^{+}>0$ 于 $U$，或者 $u^{+}\equiv 0$ 于 $U$。类似论证适用于 $u^{-}$，故式~\eqref{align:ch06-u+方程} 或~\eqref{align:ch06-u+方程} 成立。

8. 最后设 $u$ 与 $\widetilde{u}$ 是式~\eqref{align:ch06-正交} 的两个非平凡弱解。鉴于上述第 6、7 步，$\int_{U}\widetilde{u}\,\mathrm{d}x\neq 0$，故存在实常数 $\chi$ 使
\begin{align}
\int_{U}u-\chi\widetilde{u}\,\mathrm{d}x=0.
\end{align}
但由于 $u-\chi\widetilde{u}$ 也是式~\eqref{align:ch06-正交} 的弱解，第 6、7 步与式~\eqref{align:ch06-正交} 蕴含 $u\equiv\chi\widetilde{u}$ 于 $U$。因而特征值 $\lambda_{1}$ 单重。
\end{proof}

\subsection{非对称椭圆算子的特征值}\label{subsection:ch06-非对称算子特征值}

我们现考虑非散度型的一致椭圆算子：
\begin{align*}
Lu=-\sum_{i,j=1}^{n}a_{ij}u_{x_{i}x_{j}}+\sum_{i=1}^{n}b_{i}u_{x_{i}}+cu.
\end{align*}
为简单起见，设 $a_{ij}$、$b_{i}$、$c\in C^{\infty}(\overline{U})$，$U$ 开、有界、连通，$\partial U$ 光滑。另设 $a_{ij}=a_{ji}$（$i,j=1,\cdots,n$）且
\begin{align}
c\geqslant 0\quad\text{于}\ U.
\end{align}
然而注意一般算子 $L$ 不等于其形式伴随。因而我们不能如上调用\cref{section:appD-对称算子} 的抽象理论。事实上 $L$ 一般有复特征值与特征函数。

然而值得注意的是，$L$ 的主特征值为实，相应特征函数在 $U$ 内定号。

\begin{theorem}[非对称椭圆算子的主特征值]\label{theorem:ch06-非对称主特征值}
\begin{enumerate}[(i)]
\item 算子 $L$（带零边界条件）存在实特征值 $\lambda_{1}>0$，使若 $\lambda\in\mathbb{C}$ 为任一其他特征值，则 $\operatorname{Re}(\lambda)\geqslant\lambda_{1}$。
\item 存在相应特征函数 $w_{1}$，在 $U$ 内为正。
\item 特征值 $\lambda_{1}$ 单重；即若 $u$ 是式~\eqref{align:ch06-更高边界问题} 的任一解，则 $u$ 是 $w_{1}$ 的倍数。
\end{enumerate}
\end{theorem}

\begin{proof}*（初读可略）
1. 选 $m=\left[\dfrac{n}{2}\right]+3$，考虑Banach空间 $X=H^{m}(U)\cap H^{1}_{0}(U)$。由\cref{theorem:ch05-Sobolev嵌入} 的推论，$X\subset C^{2}(\overline{U})$。我们定义线性、紧算子 $A:X\to X$，令 $Af=u$，其中 $u$ 是
\begin{align}
\begin{cases}
Lu=f & \text{于}\ U,\\
u=0 & \text{于}\ \partial U
\end{cases}
\end{align}
的唯一解。其次定义锥
\begin{align*}
C=\{u\in X\mid u\geqslant 0\ \text{于}\ U\}.
\end{align*}
由极值原理，$A:C\to C$。

2. 此后固定任一函数 $w\in C$、$w\not\equiv 0$。利用强极值原理与 Hopf 引理，我们推得
\begin{align}
v>0\ \text{于}\ U,\qquad \frac{\partial v}{\partial\nu}<0\ \text{于}\ \partial U
\end{align}
其中 $v=A(w)$。记住 $w=0$ 于 $\partial U$。于是由式~\eqref{align:ch06-Aw正}, 存在常数 $\mu>0$ 使
\begin{align}
\mu v\geqslant w\quad\text{于}\ U.
\end{align}

3. 固定 $\varepsilon>0$、$\eta>0$，然后考虑方程
\begin{align}
u=\eta A[u+\varepsilon w]
\end{align}
对未知量 $u\in C$。我们断言
\begin{align}
\text{若式~\eqref{align:ch06-正交} 有解}\ u,\ \text{则}\ \eta\leqslant\mu.
\end{align}
为验证此断言，设 $u\in C$ 确实求解式~\eqref{align:ch06-正交}。由式~\eqref{align:ch06-mu关系}，计算
\begin{align*}
u\geqslant\eta A[\varepsilon w]=\eta\varepsilon v\geqslant\frac{\eta}{\mu}\varepsilon w.
\end{align*}
因而
\begin{align*}
u\geqslant\eta Au\geqslant\frac{\eta^{2}\varepsilon}{\mu}Aw=\frac{\eta^{2}\varepsilon}{\mu}v\geqslant\left(\frac{\eta}{\mu}\right)^{2}\varepsilon w.
\end{align*}
继续下去，我们推得
\begin{align*}
u\geqslant\left(\frac{\eta}{\mu}\right)^{k}\varepsilon w\qquad (k=1,\cdots),
\end{align*}
除非 $\eta\leqslant\mu$，否则矛盾。此观察确认了断言~\eqref{align:ch06-eta界}。

4. 定义
\begin{align*}
S_{\varepsilon}\triangleq\{u\in C\mid\text{存在}\ 0\leqslant\eta\leqslant 2\mu\ \text{使}\ u=\eta A[u+\varepsilon w]\}.
\end{align*}
我们断言
\begin{align}
S_{\varepsilon}\ \text{在}\ X\ \text{中无界}.
\end{align}
事实上，若非如此，则方程 $u=2\mu A[u+\varepsilon w]$ 将有解，与式~\eqref{align:ch06-eta界} 矛盾。为得出此结论，我们应用 Schaefer 不动点定理对凸集上映射的变体（\cref{theorem:ch09-2-4-Schaefer} 中陈述）。

5. 由式~\eqref{align:ch06-下解}, 存在
\begin{align}
0\leqslant\eta_{\varepsilon}\leqslant 2\mu
\end{align}
及 $v_{\varepsilon}\in C$，$\|v_{\varepsilon}\|_{X}\geqslant\dfrac{1}{\varepsilon}$，满足
\begin{align}
v_{\varepsilon}=\eta_{\varepsilon}A[v_{\varepsilon}+\varepsilon w].
\end{align}
规范化，令
\begin{align}
u_{\varepsilon}\triangleq\frac{v_{\varepsilon}}{\|v_{\varepsilon}\|_{X}}.
\end{align}
利用前述各式与算子 $A$ 的紧性，我们得到子序列 $\varepsilon_{k}\to 0$，使 $\eta_{\varepsilon_{k}}\to\eta$ 且 $u_{\varepsilon_{k}}\to u$ 于 $X$。于是式~\eqref{align:ch06-正交} 蕴含
\begin{align}
\|u\|_{X}=1,\quad u\in C.
\end{align}
由于 $u_{\varepsilon}=\eta_{\varepsilon}A\left[u_{\varepsilon}+\dfrac{\varepsilon w}{\|v_{\varepsilon}\|_{X}}\right]$，我们在极限中推得 $u=\eta Au$。由式~\eqref{align:ch06-u+方程}, $\eta>0$。因而我们可把上式改写为
\begin{align*}
\begin{cases}
Lw_{1}=\lambda_{1}w_{1} & \text{于}\ U,\\
w_{1}=0 & \text{于}\ \partial U,
\end{cases}
\end{align*}
其中 $\lambda_{1}=\eta$、$u=w_{1}$。因而 $\lambda_{1}$ 是算子 $L$（带零边界条件）的实特征值，$w_{1}\geqslant 0$ 是相应特征函数。由强极值原理与 Hopf 引理，我们有
\begin{align}
w_{1}>0\ \text{于}\ U,\qquad \frac{\partial w_{1}}{\partial\nu}<0\ \text{于}\ \partial U.
\end{align}
此外，由于\cref{section:ch06-正则性} 的正则性理论，我们知道 $w_{1}$ 光滑。

6. 上述第 1--5 步中出现的所有表达式均为实。现设 $\lambda\in\mathbb{C}$，$u$ 是
\begin{align}
\begin{cases}
Lu=\lambda u & \text{于}\ U,\\
u=0 & \text{于}\ \partial U
\end{cases}
\end{align}
的复值解。现选任一光滑函数 $w:U\to\mathbb{R}$，$w>0$ 于 $U$，令 $v\triangleq\dfrac{u}{w}$。计算
\begin{align}
\lambda v=\frac{1}{w}L(vw)\quad\text{由式~\eqref{align:ch06-更高边界问题}}
=Lv-cv-\frac{2}{w}\sum_{i,j=1}^{n}a_{ij}w_{x_{j}}v_{x_{i}}+\frac{v}{w}Lw.
\end{align}
写
\begin{align*}
Kv\triangleq-\sum_{i,j=1}^{n}a_{ij}v_{x_{i}x_{j}}+\sum_{i=1}^{n}b'_{i}v_{x_{i}},
\end{align*}
其中 $b'_{i}\triangleq b_{i}-\dfrac{2}{w}\sum_{j=1}^{n}a_{ij}w_{x_{j}}$（$i=1,\cdots,n$），我们由式~\eqref{align:ch06-lambdav} 推得
\begin{align}
Kv+\left(\frac{Lw}{w}-\lambda\right)v=0\quad\text{于}\ U.
\end{align}
取复共轭：
\begin{align}
K\overline{v}+\left(\frac{Lw}{w}-\overline{\lambda}\right)\overline{v}=0\quad\text{于}\ U.
\end{align}
其次计算
\begin{align}
K(|v|^{2})=K(v\overline{v})=\overline{v}Kv+vK\overline{v}-2\sum_{i,j=1}^{n}a_{ij}v_{x_{i}}\overline{v}_{x_{j}}\leqslant\overline{v}Kv+vK\overline{v},
\end{align}
因为对 $\xi\in\mathbb{C}^{n}$，
\begin{align*}
\sum_{i,j=1}^{n}a_{ij}\xi_{i}\overline{\xi}_{j}
=\sum_{i,j=1}^{n}a_{ij}(\operatorname{Re}(\xi_{i})\operatorname{Re}(\xi_{j})+\operatorname{Im}(\xi_{i})\operatorname{Im}(\xi_{j}))\geqslant 0.
\end{align*}
合并前述三式，我们发现
\begin{align*}
K(|v|^{2})\leqslant 2\left(\operatorname{Re}\lambda-\frac{Lw}{w}\right)|v|^{2}.
\end{align*}
现选
\begin{align}
w\triangleq w_{1}^{1-\varepsilon_{1}}\qquad (0<\varepsilon<1).
\end{align}
则
\begin{align*}
Lw=(1-\varepsilon)w_{1}^{\varepsilon}Lw_{1}+\varepsilon(1-\varepsilon)w_{1}^{1+\varepsilon}\sum_{i,j=1}^{n}a_{ij}w_{1,x_{i}}w_{1,x_{j}}+\varepsilon cw_{1}^{1-\varepsilon}\geqslant(1-\varepsilon)\lambda_{1}w.
\end{align*}
因而
\begin{align*}
K(|v|^{2})\leqslant 2(\operatorname{Re}\lambda-(1-\varepsilon)\lambda_{1})|v|^{2}\quad\text{于}\ U.
\end{align*}
于是若 $\operatorname{Re}(\lambda)\leqslant(1-\varepsilon)\lambda_{1}$，则 $K(|v|^{2})\leqslant 0$ 于 $U$。由于由式~\eqref{align:ch06-Aw正} 与~\eqref{align:ch06-Aw正} 有 $v=0$ 于 $\partial U$，我们由极值原理推得 $v=\dfrac{u}{w}=0$ 于 $U$。因而 $u\equiv 0$ 于 $U$，故 $\lambda$ 不能是特征值。此结论对每个 $\varepsilon>0$ 成立，故若 $\lambda$ 是任一复特征值，则 $\operatorname{Re}\lambda\geqslant\lambda_{1}$。

7. 最后，设 $u$ 是
\begin{align}
\begin{cases}
Lu=\lambda_{1}u & \text{于}\ U,\\
u=0 & \text{于}\ \partial U
\end{cases}
\end{align}
的任一（可能复值）解。则可证 $u$ 是 $w_{1}$ 的倍数，故 $\lambda_{1}$ 单重。
\end{proof}

\end{document}
